% [VERIFY] conclusion traces to: problem.md Insight; C02,C03,C06,C08; constraints.md (limitations); A1 open question; trace N04,N05,N10 (dead-ends); Fig 7 (interpretation only)
\section{Conclusion}
\label{sec:conclusion}
We have shown that the difficulty of training very deep convolutional networks is largely a
problem of optimization, not of representational capacity, and that reframing each block to
learn a residual with respect to its input---recovered through a parameter-free identity
shortcut---makes depth a usable resource rather than a liability. The reformulation removes
the degradation that afflicts plain networks, lets accuracy grow with depth far beyond what
was previously trainable, and does so without adding parameters, auxiliary losses, or a
specialized optimizer. The same deep representations, learned for classification, transfer
to object detection and improve localization-sensitive metrics, indicating that the depth of
representations is of central importance across visual recognition tasks rather than a
classification-specific artifact.

\subsection*{Limitations}
Several limitations bound these conclusions. First, the claim that residual learning eases
\emph{optimization} is supported empirically---by the controlled depth comparisons and by the
observation that residual layer responses are smaller in magnitude than plain ones---but is
not formally proven; the argument rests on the hypothesis that stacked nonlinear layers can
asymptotically approximate the desired residual functions, which remains an open question.
Vanishing gradients and longer training were examined and set aside as explanations: with
batch normalization the forward signals do not vanish, and training plain networks for three
times as many iterations did not close the degradation gap, so neither accounts for the
phenomenon, but ruling out these specific causes is not the same as a positive proof of the
optimization-landscape claim. Second, gains at extreme depth are not unconditional: on the
small CIFAR-10 dataset the 1202-layer network optimizes cleanly yet overfits relative to the
110-layer one, and this work applies no strong regularization to counter it. Third, the
architecture exploration is limited: a single-layer residual function offers no advantage and
is excluded, only depth (not width) is scaled, and the solver studied is stochastic gradient
descent with momentum, with no comparison to second-order methods. Finally, the headline
results depend on testing protocols that should not be mixed---10-crop and multi-scale
fully-convolutional testing report different absolute numbers, and the strongest aggregate
result is a multi-model ensemble that a single model does not reproduce.

\subsection*{Future Work}
Two concrete directions follow. The extreme-depth overfitting on small data invites combining
residual networks with explicit regularizers such as maxout or dropout, to test whether the
1202-layer model's clean optimization can be converted into a test-error improvement rather
than a regression. Separately, the residual-response analysis motivates studying stronger or
explicit regularization on the residual functions themselves and characterizing how per-layer
response magnitude scales with depth, toward a predictive rather than merely descriptive
account of why residual reformulation eases optimization.
